\documentclass[11pt,a4paper]{article}
\usepackage[margin=2.5cm]{geometry}
\usepackage{fontspec}
\usepackage{kotex}
\setmainhangulfont{Malgun Gothic}
\setsanshangulfont{Malgun Gothic}
\usepackage{amsmath,amssymb}
\usepackage{booktabs}
\usepackage{xcolor}
\usepackage{hyperref}
\usepackage{tcolorbox}
\usepackage{pgfplots}
\pgfplotsset{compat=1.18}
\tcbuselibrary{skins,breakable}

\usepackage{lscape}

\title{\Large\textbf{D-STORM 벤치마크 실험 보고서} \\[0.3em]
  \large 2026년 3월 31일 전체 통합 실측 (3회 실행 최소값)}
\author{STORM Project}
\date{}

\begin{document}
\maketitle

\section{실험 환경}
\textbf{CPU 실험 (Exp 1--7):} Apple M3 Pro 12코어, 18GB, Python 3.13.5, NumPy 2.4.3, SciPy 1.17.1, Cython 3.2.4.
\textbf{GPU + 통합 실험 (Exp 8--9):} Windows 11 Pro, Python 3.14.3, NumPy 2.4.4, SciPy 1.17.1, NVIDIA GeForce RTX 5080 (16GB), CUDA 12.9, CuPy 14.0.1, SuiteSparse:GraphBLAS 10.3.1.
STORM은 Cython C 확장(\texttt{-O3 -march=native})을 포함한 단일 구현체.
모든 수치는 3회 실행 최소값이며, 정확성은 원소별 비교로 100\% 확인.

\section{Exp 1: Facebook 소셜 네트워크}

\begin{figure}[ht]
\centering
\begin{tikzpicture}
\begin{axis}[
  ybar=2pt, width=0.85\textwidth, height=6cm, bar width=12pt,
  ylabel={실행 시간 (초)},
  symbolic x coords={NX,M-AORM,I-AORM,SciPy,STORM},
  xtick=data, ymin=0, ymax=22, enlarge x limits=0.18,
  nodes near coords={\pgfmathprintnumber[precision=1]{\pgfplotspointmeta}},
  nodes near coords style={font=\normalsize, above},
  ymajorgrids=true,
]
\addplot[fill=orange!30] coordinates {(NX,19.1)};
\addplot[fill=yellow!35] coordinates {(M-AORM,5.0)};
\addplot[fill=yellow!55] coordinates {(I-AORM,3.5)};
\addplot[fill=green!25] coordinates {(SciPy,3.0)};
\addplot[fill=blue!60] coordinates {(STORM,1.2)};
\end{axis}
\end{tikzpicture}
\caption{Facebook ($n\!=\!4{,}039$, $m\!=\!176{,}468$, $d\!=\!8$) APSP}
\label{fig:fb}
\end{figure}

\begin{table}[ht]\centering
\caption{Facebook APSP 결과}
\begin{tabular}{@{}lrrr@{}}
\toprule
\textbf{방법} & \textbf{시간 (s)} & \textbf{STORM 대비} & \textbf{NX 대비} \\
\midrule
\textbf{STORM} & \textbf{1.156} & 1.0$\times$ & 16.5$\times$ \\
SciPy (C) & 2.960 & 0.39$\times$ & 6.5$\times$ \\
I-AORM & 3.502 & 0.33$\times$ & 5.5$\times$ \\
M-AORM & 4.970 & 0.23$\times$ & 3.8$\times$ \\
NetworkX & 19.110 & 0.06$\times$ & 1.0$\times$ \\
\bottomrule
\end{tabular}
\end{table}

\textbf{발견 1:} STORM은 I-AORM 대비 \textbf{3.0배}, C로 구현된 SciPy Dijkstra 대비 \textbf{2.6배}, NetworkX 대비 \textbf{16.5배} 빠르다. M-AORM 대비로는 4.3배이다.

\section{Exp 2: BA 그래프 확장성}

\begin{figure}[ht]
\centering
\begin{tikzpicture}
\begin{axis}[
  width=0.85\textwidth, height=6cm,
  xlabel={노드 수}, ylabel={실행 시간 (초)},
  xmode=log, ymode=log, log basis x=10, log basis y=10,
  legend style={at={(0.02,0.98)}, anchor=north west, font=\small},
  grid=both, mark size=3pt,
]
\addplot[orange,thick,mark=triangle*] coordinates {(500,0.0702)(1000,0.3004)(2000,1.2752)(3000,3.0843)(5000,8.7188)};
\addlegendentry{NetworkX}
\addplot[yellow!70!black,thick,mark=+,mark size=3.5pt] coordinates {(500,0.0177)(1000,0.0885)(2000,0.3224)(3000,0.6690)};
\addlegendentry{I-AORM}
\addplot[green!60!black,thick,mark=pentagon*] coordinates {(500,0.0283)(1000,0.1165)(2000,0.4792)(3000,1.0982)(5000,3.1485)};
\addlegendentry{SciPy (C)}
\addplot[blue,very thick,mark=square*,mark options={fill=blue!50,scale=1.2}] coordinates {(500,0.0174)(1000,0.0686)(2000,0.2817)(3000,0.6489)(5000,1.8723)};
\addlegendentry{\textbf{STORM}}
\end{axis}
\end{tikzpicture}
\caption{BA 그래프 확장성 (로그-로그)}
\label{fig:ba}
\end{figure}

\textbf{발견 2:} STORM이 모든 스케일에서 최저 시간. SciPy(C) 대비 $n=500$에서 1.6배, $n=5{,}000$에서 1.7배. I-AORM은 $n > 3{,}000$에서 측정 불가.

\section{Exp 3: 위상별 성능 — 지름 효과}

\begin{table}[ht]\centering
\caption{위상별 비교 ($n \approx 2{,}000$, 초)}
\begin{tabular}{@{}lrrrrrr@{}}
\toprule
\textbf{위상} & $d$ & \textbf{STORM} & \textbf{I-AORM} & \textbf{SciPy} & \textbf{NX} & \textbf{vs SciPy} \\
\midrule
BA & 5 & \textbf{0.283} & 0.324 & 0.480 & 1.235 & 1.7$\times$ \\
ER & 6 & \textbf{0.320} & 0.379 & 0.488 & 1.361 & 1.5$\times$ \\
WS & 13 & \textbf{0.291} & 0.582 & 0.422 & 1.236 & 1.5$\times$ \\
Grid & 88 & 0.475 & 3.203 & \textbf{0.303} & 1.032 & 0.6$\times$ \\
PLC & 5 & \textbf{0.284} & 0.324 & 0.481 & 1.256 & 1.7$\times$ \\
\bottomrule
\end{tabular}
\end{table}

\begin{figure}[ht]
\centering
\begin{tikzpicture}
\begin{axis}[
  width=0.85\textwidth, height=5.5cm,
  xlabel={그래프 지름 $d$}, ylabel={STORM / SciPy 속도비},
  xmin=0, xmax=100, ymin=0, ymax=2.0,
  grid=major, mark size=4pt,
]
\addplot[blue,very thick,mark=square*,only marks,nodes near coords,
  point meta=explicit symbolic,
  every node near coord/.style={anchor=south west, font=\small}]
  coordinates {(5,1.70)[BA] (6,1.53)[ER] (13,1.45)[WS] (88,0.64)[Grid] (5,1.69)[PLC]};
\draw[red,thick,dashed] (axis cs:0,1) -- (axis cs:100,1);
\fill[green!10] (axis cs:0,1) rectangle (axis cs:20,2);
\node[green!50!black,font=\small\bfseries] at (axis cs:10,1.85) {STORM 우위};
\node[red!50!black,font=\small\bfseries] at (axis cs:60,0.4) {SciPy 우위};
\end{axis}
\end{tikzpicture}
\caption{그래프 지름과 STORM/SciPy 속도비}
\label{fig:diam}
\end{figure}

\textbf{발견 3:} \textbf{그래프 지름이 점진적 APSP의 성능 결정 요인이다.} $d \leq 13$에서 STORM이 SciPy 대비 1.5--1.7배 우위. Grid($d=88$)에서는 SciPy가 1.6배 빠르지만, Cython 커널 덕분에 순수 Python(8.7배 느림)에서 크게 개선됨. STORM은 I-AORM을 \textbf{모든 위상에서} 능가한다.

\section{Exp 4: 홉 제약 APSP}

\begin{figure}[ht]
\centering
\begin{tikzpicture}
\begin{axis}[
  width=0.85\textwidth, height=6cm,
  xlabel={홉 제약 $k$}, ylabel={실행 시간 (초)},
  xmin=1.5, xmax=8.5, xtick={2,...,8}, ymin=0, ymax=21,
  legend style={at={(0.02,0.98)}, anchor=north west, font=\small},
  grid=major, mark size=3pt,
]
\addplot[orange,thick,mark=triangle*] coordinates {(2,0.764)(3,4.541)(4,10.106)(5,17.520)(6,18.419)(7,18.867)(8,19.162)};
\addlegendentry{NetworkX}
\addplot[yellow!70!black,thick,mark=+,mark size=3.5pt] coordinates {(2,0.545)(3,0.980)(4,1.397)(5,1.858)(6,2.301)(7,2.723)(8,3.206)};
\addlegendentry{I-AORM}
\addplot[blue,very thick,mark=square*] coordinates {(2,0.080)(3,0.301)(4,0.617)(5,0.959)(6,1.091)(7,1.145)(8,1.178)};
\addlegendentry{\textbf{STORM}}
\draw[<->,thick,green!50!black] (axis cs:2,0.080) -- (axis cs:2,0.545) node[midway,right,font=\small\bfseries] {6.8$\times$};
\draw[<->,thick,green!50!black] (axis cs:5,0.959) -- (axis cs:5,17.520) node[midway,right,font=\small\bfseries] {18.3$\times$};
\end{axis}
\end{tikzpicture}
\caption{Facebook 홉 제약 APSP}
\label{fig:hop}
\end{figure}

\textbf{발견 4:} STORM의 최대 강점. $k=2$에서 I-AORM 대비 \textbf{6.8배} (희소 $\mathbf{R}^{(2)*}$에서 Cython 커널 효과 극대화). $k=5$에서 NX 대비 \textbf{18.3배} (정확도 89.4\%). 교통 혼잡 예측($k=3$--$5$), 감염 추적($k=2$--$4$)에 직접 적용 가능.

\section{Exp 5: 점진적 희소성과 수렴}

\begin{figure}[ht]
\centering
\begin{tikzpicture}
\begin{axis}[
  width=0.85\textwidth, height=5.5cm,
  xlabel={차수 $k$}, ylabel={nnz($\mathbf{R}^{(k)*}$) / 백만},
  axis y line*=left, ymin=0, ymax=7, xmin=1, xmax=8, xtick={1,...,8},
  ylabel style={blue}, yticklabel style={blue},
  legend style={at={(0.02,0.98)}, anchor=north west, font=\small},
]
\addplot[blue,very thick,mark=square*,fill=blue!30] coordinates {(1,0.176)(2,2.716)(3,3.982)(4,5.862)(5,2.565)(6,0.677)(7,0.315)(8,0.016)};
\addlegendentry{nnz($\mathbf{R}^{(k)*}$)}
\end{axis}
\begin{axis}[
  width=0.85\textwidth, height=5.5cm,
  ylabel={정확도 $\mathcal{C}$}, axis y line*=right, axis x line=none,
  ymin=0, ymax=1.15, xmin=1, xmax=8,
  ylabel style={red!70!black}, yticklabel style={red!70!black},
  legend style={at={(0.55,0.45)}, anchor=north west, font=\small},
]
\addplot[red!70!black,thick,dashed,mark=triangle*] coordinates {(1,0.003)(2,0.093)(3,0.291)(4,0.681)(5,0.894)(6,0.961)(7,0.998)(8,1.0)};
\addlegendentry{정확도}
\end{axis}
\end{tikzpicture}
\caption{Facebook 점진적 희소성 (이중 축)}
\label{fig:sparse}
\end{figure}

\textbf{발견 5:} nnz($\mathbf{R}^{(k)*}$)가 $k=4$에서 590만으로 정점 후 $k=8$에서 1.6만(370배 감소). STORM은 이 희소성을 자동으로 활용하여 후반 반복을 가속. $k=5$에서 정확도 89\%, $k=7$에서 99.8\% --- 5회 반복만으로 대부분의 정보 획득.

\section{Exp 6: D-STORM 동적 갱신}

\begin{table}[ht]\centering
\caption{D-STORM 엣지 추가 (30회 평균)}
\begin{tabular}{@{}rrrrrr@{}}
\toprule
$n$ & \textbf{추가 (ms)} & \textbf{영향 쌍} & \textbf{전체 (ms)} & \textbf{배수} \\
\midrule
500 & 0.80 & 20 & 17.2 & \textbf{22$\times$} \\
1,000 & 2.32 & 42 & 68.9 & \textbf{30$\times$} \\
2,000 & 8.52 & 92 & 283.7 & \textbf{33$\times$} \\
\bottomrule
\end{tabular}
\end{table}

\textbf{발견 6:} D-STORM 엣지 추가가 전체 재계산 대비 22--33배 빠름. 그래프가 커질수록 speedup 증가(500: 22배 → 2000: 33배) --- 영향 비율이 $n^2$ 대비 감소하기 때문.

\section{Exp 7: GraphBLAS 베이스라인 비교}

SuiteSparse:GraphBLAS v10.3.1 (C 라이브러리, CFFI 바인딩)을 사용하여 동일 알고리즘의 GraphBLAS 구현과 D-STORM을 직접 비교한다. Python 3.14, Windows 11 환경에서 측정. 3회 실행 최소값.

\subsection{GraphBLAS 구현 방법}
\begin{itemize}
  \item \textbf{GB-bfs}: 소스별 masked BFS (\texttt{GrB\_vxm} + 구조적 보수 마스크). LAGraph 표준 패턴.
  \item \textbf{GB-frontier}: D-STORM의 cumulative footprint 알고리즘을 GraphBLAS \texttt{GrB\_mxm} (LOR\_LAND 세미링) + 보수 마스크로 재구현. 동일 알고리즘, GraphBLAS C 커널.
\end{itemize}

\subsection{Facebook 5-방법 비교 (GraphBLAS 포함)}

\begin{table}[ht]\centering
\caption{Facebook APSP: GraphBLAS 포함 전체 비교 ($n\!=\!4{,}039$)}
\begin{tabular}{@{}lrrrr@{}}
\toprule
\textbf{방법} & \textbf{시간 (s)} & \textbf{NX 대비} & \textbf{SciPy 대비} & \textbf{커널} \\
\midrule
\textbf{D-STORM} & \textbf{1.548} & 8.8$\times$ & 1.2$\times$ & Cython C \\
SciPy (C Dijkstra) & 1.841 & 7.4$\times$ & 1.0$\times$ & C \\
GB-frontier & 1.865 & 7.3$\times$ & 1.0$\times$ & GraphBLAS C \\
GB-bfs & 4.041 & 3.4$\times$ & 0.5$\times$ & GraphBLAS C \\
NetworkX & 13.669 & 1.0$\times$ & --- & Python \\
\bottomrule
\end{tabular}
\end{table}

\begin{figure}[ht]
\centering
\begin{tikzpicture}
\begin{axis}[
  ybar=2pt, width=0.85\textwidth, height=6cm, bar width=11pt,
  ylabel={실행 시간 (초)},
  symbolic x coords={NX,GB-bfs,SciPy,GB-frontier,D-STORM},
  xtick=data, ymin=0, ymax=16, enlarge x limits=0.15,
  nodes near coords={\pgfmathprintnumber[precision=2]{\pgfplotspointmeta}},
  nodes near coords style={font=\small, above},
  ymajorgrids=true,
]
\addplot[fill=orange!30] coordinates {(NX,13.669)};
\addplot[fill=red!25] coordinates {(GB-bfs,4.041)};
\addplot[fill=green!25] coordinates {(SciPy,1.841)};
\addplot[fill=cyan!35] coordinates {(GB-frontier,1.865)};
\addplot[fill=blue!60] coordinates {(D-STORM,1.548)};
\end{axis}
\end{tikzpicture}
\caption{Facebook APSP: GraphBLAS 포함 비교. GB-frontier(동일 알고리즘, GraphBLAS C 커널)와 SciPy가 거의 동등하며, D-STORM이 1.2배 빠르다.}
\label{fig:fb-gb}
\end{figure}

\textbf{발견 7a:} GB-frontier(동일 알고리즘, GraphBLAS C 커널)는 SciPy와 거의 동등한 성능(1.865s vs 1.841s). D-STORM은 GB-frontier보다 \textbf{1.2배} 빠르다 --- Cython fused pruning 커널이 GraphBLAS의 범용 SpMM보다 이 워크로드에서 효율적임을 시사한다. GB-bfs(소스별 BFS)는 Python-level 루프 오버헤드로 인해 4.0s로 가장 느리다.

\subsection{BA 그래프 확장성 (GraphBLAS 포함)}

\begin{table}[ht]\centering
\caption{BA 그래프 확장성: GraphBLAS 포함 (초)}
\begin{tabular}{@{}rrrrrr@{}}
\toprule
$n$ & \textbf{D-STORM} & \textbf{GB-frontier} & \textbf{SciPy} & \textbf{GB-bfs} & \textbf{STORM/GB-f} \\
\midrule
500 & \textbf{0.017} & 0.034 & 0.016 & 0.103 & 2.0$\times$ \\
1,000 & \textbf{0.081} & 0.121 & 0.069 & 0.319 & 1.5$\times$ \\
2,000 & \textbf{0.328} & 0.467 & 0.294 & 1.108 & 1.4$\times$ \\
3,000 & \textbf{0.760} & 1.087 & 0.691 & 2.503 & 1.4$\times$ \\
\bottomrule
\end{tabular}
\end{table}

\begin{figure}[ht]
\centering
\begin{tikzpicture}
\begin{axis}[
  width=0.85\textwidth, height=6cm,
  xlabel={노드 수}, ylabel={실행 시간 (초)},
  xmode=log, ymode=log, log basis x=10, log basis y=10,
  legend style={at={(0.02,0.98)}, anchor=north west, font=\small},
  grid=both, mark size=3pt,
]
\addplot[red,thick,mark=diamond*] coordinates {(500,0.103)(1000,0.319)(2000,1.108)(3000,2.503)};
\addlegendentry{GB-bfs}
\addplot[cyan!70!black,thick,mark=o] coordinates {(500,0.034)(1000,0.121)(2000,0.467)(3000,1.087)};
\addlegendentry{GB-frontier}
\addplot[green!60!black,thick,mark=pentagon*] coordinates {(500,0.016)(1000,0.069)(2000,0.294)(3000,0.691)};
\addlegendentry{SciPy (C)}
\addplot[blue,very thick,mark=square*,mark options={fill=blue!50,scale=1.2}] coordinates {(500,0.017)(1000,0.081)(2000,0.328)(3000,0.760)};
\addlegendentry{\textbf{D-STORM}}
\end{axis}
\end{tikzpicture}
\caption{BA 그래프 확장성: GraphBLAS 포함 (로그--로그)}
\label{fig:ba-gb}
\end{figure}

\textbf{발견 7b:} D-STORM은 모든 스케일에서 GB-frontier보다 \textbf{1.4--2.0배} 빠르다. 이 차이는 그래프 크기에 관계없이 안정적이다. 한편, SciPy(C)가 BA 그래프에서 D-STORM과 비슷하거나 더 빠른 경우도 있음에 주의 --- 소규모 BA 그래프에서는 C Dijkstra의 상수 계수가 작기 때문이다.

\subsection{위상별 비교 (GraphBLAS 포함)}

\begin{table}[ht]\centering
\caption{위상별 비교: GraphBLAS 포함 ($n \approx 2{,}000$, 초)}
\begin{tabular}{@{}lrrrrr@{}}
\toprule
\textbf{위상} & $d$ & \textbf{D-STORM} & \textbf{GB-frontier} & \textbf{SciPy} & \textbf{STORM/GB-f} \\
\midrule
BA & 5 & \textbf{0.327} & 0.459 & 0.292 & 1.4$\times$ \\
ER & 6 & \textbf{0.380} & 0.468 & 0.298 & 1.2$\times$ \\
WS & 13 & \textbf{0.379} & 0.484 & 0.276 & 1.3$\times$ \\
Grid & 88 & 3.209 & \textbf{0.573} & \underline{0.149} & 0.18$\times$ \\
PLC & 5 & \textbf{0.326} & 0.457 & 0.291 & 1.4$\times$ \\
\bottomrule
\end{tabular}
\end{table}

\textbf{발견 7c:} 소규모 지름 그래프($d \leq 13$)에서 D-STORM이 GB-frontier보다 1.2--1.4배 빠르다. 그러나 \textbf{Grid($d=88$)에서는 역전}: GB-frontier가 D-STORM보다 \textbf{5.6배} 빠르다. 이는 Grid에서 D-STORM의 Python/Cython 루프가 88회 반복되면서 오버헤드가 누적되는 반면, GraphBLAS의 C SpMM 커널은 반복당 비용이 더 낮기 때문이다. SciPy(C Dijkstra)는 Grid에서 여전히 최고 성능(0.149s).

\subsection{해석과 시사점}

\begin{tcolorbox}[colback=blue!5, colframe=blue!50!black, title=GraphBLAS 비교 핵심 결론]
\begin{enumerate}
  \item \textbf{동일 알고리즘에서 D-STORM이 GB-frontier보다 1.2--1.4배 빠르다} (소세계 그래프). Cython fused pruning 커널이 GraphBLAS 범용 SpMM보다 이 워크로드에 특화되어 있기 때문.
  \item \textbf{고지름 그래프에서 GraphBLAS가 D-STORM을 크게 능가한다} (Grid: 5.6배). 반복 횟수가 많아지면 GraphBLAS의 C-level 반복 효율이 결정적.
  \item \textbf{GB-bfs(소스별 BFS)는 Python 루프 오버헤드로 인해 비경쟁적}. C/C++로 전체 소스 루프까지 구현한 LAGraph 수준의 비교가 추가 필요.
  \item \textbf{D-STORM의 조건부 우위가 확인}: 소세계 그래프(소셜, PPI 등)에서는 D-STORM이 GraphBLAS 재구현보다 빠르며, 이는 fused pruning + footprint 갱신의 단일 패스 설계 덕분.
\end{enumerate}
\end{tcolorbox}

\section{Exp 8: GPU 가속 (CUDA / CuPy)}

NVIDIA GeForce RTX 5080 (16GB, CUDA 12.9), CuPy 14.0.1 환경에서 D-STORM의 GPU 가속 버전 3종을 구현하고 측정한다. 1회 warmup 후 3회 실행 최소값.

\subsection{GPU 실행 전략}

\begin{table}[ht]\centering
\caption{GPU 실행 전략 비교}
\small
\begin{tabular}{@{}llll@{}}
\toprule
\textbf{전략} & \textbf{SpMM 커널} & \textbf{Prune 커널} & \textbf{최적 대상} \\
\midrule
GPU-Dense & cuBLAS dense matmul & GPU element-wise & 소규모 ($n < 5$K) \\
GPU-Sparse & cuSPARSE SpMM & GPU sparse ops & 대규모 희소 그래프 \\
\bottomrule
\end{tabular}
\end{table}

\subsection{Facebook 전체 비교 (GPU 포함)}

\begin{table}[ht]\centering
\caption{Facebook APSP: GPU 포함 전체 비교 ($n\!=\!4{,}039$)}
\begin{tabular}{@{}lrrrrl@{}}
\toprule
\textbf{방법} & \textbf{시간 (s)} & \textbf{vs CPU} & \textbf{vs NX} & \textbf{vs SciPy} & \textbf{실행 위치} \\
\midrule
\textbf{GPU-Dense} & \textbf{0.148} & \textbf{10.1$\times$} & 85.5$\times$ & 12.3$\times$ & RTX 5080 \\
GPU-Sparse & 0.161 & 9.3$\times$ & 78.9$\times$ & 11.4$\times$ & RTX 5080 \\
CPU D-STORM & 1.492 & 1.0$\times$ & 8.5$\times$ & 1.2$\times$ & CPU (Cython) \\
SciPy (C) & 1.827 & 0.8$\times$ & 6.9$\times$ & 1.0$\times$ & CPU (C) \\
NetworkX & 12.677 & 0.1$\times$ & 1.0$\times$ & --- & CPU (Python) \\
\bottomrule
\end{tabular}
\end{table}

\begin{figure}[ht]
\centering
\begin{tikzpicture}
\begin{axis}[
  ybar=2pt, width=0.85\textwidth, height=6.5cm, bar width=10pt,
  ylabel={실행 시간 (초)},
  symbolic x coords={NX,SciPy,CPU-STORM,GPU-Sparse,GPU-Dense},
  xtick=data, ymin=0, ymax=14, enlarge x limits=0.12,
  x tick label style={rotate=25, anchor=east, font=\small},
  nodes near coords={\pgfmathprintnumber[precision=2]{\pgfplotspointmeta}},
  nodes near coords style={font=\small, above},
  ymajorgrids=true,
]
\addplot[fill=orange!30] coordinates {(NX,12.677)};
\addplot[fill=green!25] coordinates {(SciPy,1.827)};
\addplot[fill=blue!40] coordinates {(CPU-STORM,1.492)};
\addplot[fill=purple!50] coordinates {(GPU-Sparse,0.161)};
\addplot[fill=purple!70] coordinates {(GPU-Dense,0.148)};
\end{axis}
\end{tikzpicture}
\caption{Facebook APSP: GPU vs CPU. GPU-Dense가 CPU D-STORM 대비 10.1배, NetworkX 대비 85.5배 빠르다.}
\label{fig:fb-gpu}
\end{figure}

\textbf{발견 8a:} GPU-Dense(cuBLAS)가 Facebook에서 최고 성능. CPU D-STORM 대비 \textbf{10.1배}, SciPy(C) 대비 \textbf{12.3배}, NetworkX 대비 \textbf{85.5배} 빠르다. 정확성은 CPU 결과와 원소별 100\% 일치(PASS).

\subsection{BA 그래프 확장성 (GPU 포함)}

\begin{table}[ht]\centering
\caption{BA 그래프 확장성: GPU vs CPU (초, 3회 최소값)}
\begin{tabular}{@{}rrrrrrr@{}}
\toprule
$n$ & \textbf{GPU-Dense} & \textbf{GPU-Sparse} & \textbf{CPU-STORM} & \textbf{SciPy} & \textbf{GPU speedup} \\
\midrule
500 & \textbf{0.002} & 0.010 & 0.015 & 0.016 & 9.9$\times$ \\
1,000 & \textbf{0.006} & 0.013 & 0.079 & 0.069 & 12.4$\times$ \\
2,000 & 0.028 & \textbf{0.022} & 0.323 & 0.292 & 14.5$\times$ \\
3,000 & 0.068 & \textbf{0.039} & 0.732 & 0.677 & 18.7$\times$ \\
5,000 & 0.205 & \textbf{0.097} & 2.187 & 1.966 & 22.7$\times$ \\
\bottomrule
\end{tabular}
\end{table}

\begin{figure}[ht]
\centering
\begin{tikzpicture}
\begin{axis}[
  width=0.85\textwidth, height=6.5cm,
  xlabel={노드 수}, ylabel={실행 시간 (초)},
  xmode=log, ymode=log, log basis x=10, log basis y=10,
  legend style={at={(0.02,0.98)}, anchor=north west, font=\small},
  grid=both, mark size=3pt,
]
\addplot[green!60!black,thick,mark=pentagon*] coordinates {(500,0.016)(1000,0.069)(2000,0.292)(3000,0.677)(5000,1.966)};
\addlegendentry{SciPy (C)}
\addplot[blue,thick,mark=square*] coordinates {(500,0.015)(1000,0.079)(2000,0.323)(3000,0.732)(5000,2.187)};
\addlegendentry{CPU D-STORM}
\addplot[purple!80!black,very thick,mark=diamond*,mark options={fill=purple!40,scale=1.2}] coordinates {(500,0.002)(1000,0.006)(2000,0.022)(3000,0.039)(5000,0.097)};
\addlegendentry{\textbf{GPU (best)}}
\end{axis}
\end{tikzpicture}
\caption{BA 확장성: GPU(최고) vs CPU. GPU speedup은 그래프 크기에 비례하여 증가한다.}
\label{fig:ba-gpu}
\end{figure}

\textbf{발견 8b:} GPU speedup이 그래프 크기에 비례하여 증가: $n\!=\!500$에서 \textbf{9.9배} → $n\!=\!5{,}000$에서 \textbf{22.7배}. $n \leq 1{,}000$에서는 GPU-Dense(cuBLAS)가 최고, $n \geq 2{,}000$에서는 GPU-Sparse(cuSPARSE)가 최고. 이는 소규모에서 cuBLAS의 $O(n^3)$ matmul이 높은 FP throughput으로 빠르고, 대규모에서 cuSPARSE의 SpMM이 희소성을 활용하기 때문이다.

\subsection{GPU 전략 선택 가이드}

\begin{tcolorbox}[colback=purple!5, colframe=purple!50!black, title=GPU 전략 선택]
\begin{itemize}
  \item $n < 1{,}000$: \textbf{GPU-Dense} (cuBLAS matmul) --- 가장 빠름. $O(n^2)$ 메모리 허용 범위.
  \item $1{,}000 \leq n \leq 10{,}000$: \textbf{GPU-Sparse} (cuSPARSE SpMM) --- 희소성 활용, 메모리 효율적.
  \item $n > 10{,}000$: \textbf{GPU-Sparse} 전용 --- dense 행렬이 VRAM 초과 ($n^2 \times 4$ bytes $> 400$MB).
\end{itemize}
\end{tcolorbox}

\subsection{해석과 시사점}

\begin{tcolorbox}[colback=purple!5, colframe=purple!50!black, title=GPU 가속 핵심 결론]
\begin{enumerate}
  \item \textbf{D-STORM의 frontier 알고리즘은 GPU 병렬화에 적합하다.} cuSPARSE SpMM이 frontier 확장을 벡터화하고, cuBLAS가 소규모 dense 연산을 가속한다.
  \item \textbf{GPU speedup이 그래프 크기에 비례하여 증가한다} ($9.9\times$ → $22.7\times$). 이는 SpMM의 병렬 이득이 커지고 Python 오버헤드 비율이 줄어들기 때문.
  \item \textbf{CPU D-STORM의 Cython 최적화는 여전히 가치 있다.} GPU 없는 환경에서도 SciPy(C)와 경쟁 가능.
\end{enumerate}
\end{tcolorbox}

\section{Exp 9: 전체 통합 벤치마크 (11 methods $\times$ 11 datasets)}

모든 기법의 세부 변종 11종을 11개 데이터셋에서 동일 프로토콜(warmup 1회 + 3회 실행 최소값)로 측정. Windows 11, Python 3.14, RTX 5080, CUDA 12.9 환경. 총 121개 측정점 중 110개 유효(중복 제거).

\subsection{측정 대상 기법 (11종)}

\begin{table}[ht]\centering
\caption{전체 측정 대상 기법과 세부 분류}
\small
\begin{tabular}{@{}llll@{}}
\toprule
\textbf{기법} & \textbf{계열} & \textbf{핵심 커널} & \textbf{실행 위치} \\
\midrule
NetworkX & 참조 & Python BFS & CPU \\
SciPy & 참조 & C Dijkstra & CPU \\
GraphBLAS-bfs & GraphBLAS & \texttt{GrB\_vxm} 소스별 BFS & CPU (GraphBLAS C) \\
GraphBLAS-frontier & GraphBLAS & \texttt{GrB\_mxm} (LOR\_LAND) & CPU (GraphBLAS C) \\
I-AORM & AORM & Dense Python loop & CPU \\
M-AORM & AORM & Dense BLAS matmul & CPU \\
D-STORM-Sparse & D-STORM & Cython fused SpMM + prune & CPU \\
D-STORM-Dense & D-STORM & Dense BLAS matmul & CPU \\
GPU-Dense & D-STORM GPU & cuBLAS dense matmul & RTX 5080 \\
GPU-Sparse & D-STORM GPU & cuSPARSE SpMM & RTX 5080 \\
\bottomrule
\end{tabular}
\end{table}

\subsection{Exp 9a: Facebook 전체 비교 (11 methods)}

\begin{table}[ht]\centering
\caption{Facebook ($n\!=\!4{,}039$, $m\!=\!176{,}468$, $d\!=\!8$): 전체 11종 비교}
\small
\begin{tabular}{@{}rlrrrr@{}}
\toprule
\textbf{순위} & \textbf{기법} & \textbf{시간 (s)} & \textbf{vs 1위} & \textbf{vs SciPy} & \textbf{vs NX} \\
\midrule
1 & \textbf{GPU-Dense} & \textbf{0.126} & 1.0$\times$ & 14.5$\times$ & 101.3$\times$ \\
2 & GPU-Sparse & 0.176 & 1.4$\times$ & 10.4$\times$ & 72.5$\times$ \\
4 & D-STORM-Sparse & 1.506 & 12.0$\times$ & 1.2$\times$ & 8.5$\times$ \\
5 & D-STORM-Dense & 1.718 & 13.7$\times$ & 1.1$\times$ & 7.4$\times$ \\
6 & SciPy (C) & 1.821 & 14.5$\times$ & 1.0$\times$ & 7.0$\times$ \\
7 & GB-frontier & 1.991 & 15.8$\times$ & 0.9$\times$ & 6.4$\times$ \\
8 & M-AORM & 3.230 & 25.7$\times$ & 0.6$\times$ & 3.9$\times$ \\
9 & GB-bfs & 4.302 & 34.2$\times$ & 0.4$\times$ & 3.0$\times$ \\
10 & I-AORM & 7.915 & 62.9$\times$ & 0.2$\times$ & 1.6$\times$ \\
11 & NetworkX & 12.747 & 101.3$\times$ & 0.1$\times$ & 1.0$\times$ \\
\bottomrule
\end{tabular}
\end{table}

\begin{figure}[ht]
\centering
\begin{tikzpicture}
\begin{axis}[
  xbar, width=0.88\textwidth, height=8cm, bar width=6pt,
  xlabel={실행 시간 (초)},
  symbolic y coords={NetworkX,I-AORM,GB-bfs,M-AORM,GB-frontier,SciPy,
    D-STORM-Dense,D-STORM-Sparse,GPU-Sparse,GPU-Dense},
  ytick=data, xmin=0, xmax=14, enlarge y limits=0.06,
  nodes near coords={\pgfmathprintnumber[precision=2]{\pgfplotspointmeta}},
  nodes near coords style={font=\scriptsize, anchor=west},
  y tick label style={font=\small},
]
\addplot[fill=orange!30] coordinates {(12.747,NetworkX)};
\addplot[fill=orange!20] coordinates {(7.915,I-AORM)};
\addplot[fill=red!20] coordinates {(4.302,GB-bfs)};
\addplot[fill=yellow!30] coordinates {(3.230,M-AORM)};
\addplot[fill=cyan!25] coordinates {(1.991,GB-frontier)};
\addplot[fill=green!25] coordinates {(1.821,SciPy)};
\addplot[fill=blue!25] coordinates {(1.718,D-STORM-Dense)};
\addplot[fill=blue!40] coordinates {(1.506,D-STORM-Sparse)};
\addplot[fill=purple!40] coordinates {(0.176,GPU-Sparse)};
\addplot[fill=purple!60] coordinates {(0.126,GPU-Dense)};
\end{axis}
\end{tikzpicture}
\caption{Facebook APSP 전체 11종 비교 (수평 막대, 짧을수록 빠름)}
\label{fig:fb-all}
\end{figure}

\subsection{Exp 9b: BA 확장성 (11 methods $\times$ 5 scales)}

\begin{table}[ht]\centering
\caption{BA 그래프 확장성: 전체 11종 (초)}
\scriptsize
\resizebox{\textwidth}{!}{%
\begin{tabular}{@{}lrrrrr@{}}
\toprule
\textbf{기법} & $n\!=\!500$ & $n\!=\!1{,}000$ & $n\!=\!2{,}000$ & $n\!=\!3{,}000$ & $n\!=\!5{,}000$ \\
\midrule
GPU-Dense & \textbf{0.001} & \textbf{0.006} & 0.030 & 0.068 & 0.207 \\
GPU-Sparse & 0.011 & 0.015 & \textbf{0.031} & \textbf{0.042} & \textbf{0.094} \\
D-STORM-Dense & 0.010 & 0.062 & 0.329 & 0.646 & 2.061 \\
D-STORM-Sparse & 0.016 & 0.080 & 0.395 & 0.757 & 2.223 \\
SciPy (C) & 0.017 & 0.072 & 0.300 & 0.699 & 2.004 \\
I-AORM & 0.021 & 0.098 & 0.459 & 0.998 & 2.802 \\
M-AORM & 0.018 & 0.095 & 0.483 & 1.102 & 3.427 \\
GB-frontier & 0.034 & 0.123 & 0.476 & 1.075 & 2.950 \\
GB-bfs & 0.109 & 0.343 & 1.162 & 2.547 & 6.718 \\
NetworkX & 0.047 & 0.228 & 0.967 & 2.419 & 6.720 \\
\bottomrule
\end{tabular}}
\end{table}

\begin{figure}[ht]
\centering
\begin{tikzpicture}
\begin{axis}[
  width=0.88\textwidth, height=7cm,
  xlabel={노드 수 $n$}, ylabel={실행 시간 (초)},
  xmode=log, ymode=log, log basis x=10, log basis y=10,
  legend style={at={(0.02,0.98)}, anchor=north west, font=\scriptsize, legend columns=2},
  grid=both, mark size=2.5pt,
]
\addplot[orange,thick,mark=triangle*] coordinates {(500,0.047)(1000,0.228)(2000,0.967)(3000,2.419)(5000,6.720)};
\addlegendentry{NetworkX}
\addplot[orange!60,thick,mark=x,mark size=3pt] coordinates {(500,0.021)(1000,0.098)(2000,0.459)(3000,0.998)(5000,2.802)};
\addlegendentry{I-AORM}
\addplot[yellow!60!black,thick,mark=+,mark size=3pt] coordinates {(500,0.018)(1000,0.095)(2000,0.483)(3000,1.102)(5000,3.427)};
\addlegendentry{M-AORM}
\addplot[green!60!black,thick,mark=pentagon*] coordinates {(500,0.017)(1000,0.072)(2000,0.300)(3000,0.699)(5000,2.004)};
\addlegendentry{SciPy}
\addplot[cyan!70!black,thick,mark=o] coordinates {(500,0.034)(1000,0.123)(2000,0.476)(3000,1.075)(5000,2.950)};
\addlegendentry{GB-frontier}
\addplot[blue,thick,mark=square*] coordinates {(500,0.016)(1000,0.080)(2000,0.395)(3000,0.757)(5000,2.223)};
\addlegendentry{D-STORM-Sparse}
\addplot[purple!80!black,very thick,mark=diamond*,mark options={fill=purple!40,scale=1.3}] coordinates {(500,0.001)(1000,0.006)(2000,0.031)(3000,0.042)(5000,0.094)};
\addlegendentry{\textbf{GPU (best)}}
\end{axis}
\end{tikzpicture}
\caption{BA 확장성: 주요 기법 (로그--로그). GPU가 모든 스케일에서 압도적.}
\label{fig:ba-all}
\end{figure}

\textbf{발견 9b:} GPU(best)는 $n\!=\!500$에서 CPU 최고(GPU-Dense 0.001s vs SciPy 0.017s = \textbf{17배})부터 $n\!=\!5{,}000$에서(GPU-Sparse 0.094s vs SciPy 2.004s = \textbf{21배})까지 일관된 우위. CPU 내에서는 D-STORM-Dense $\approx$ SciPy $>$ D-STORM-Sparse $>$ I-AORM $>$ M-AORM $>$ GB-frontier $>$ NetworkX $\approx$ GB-bfs.

\subsection{Exp 9c: 위상별 비교 (11 methods $\times$ 5 topologies)}

\begin{table}[ht]\centering
\caption{위상별 비교 ($n \approx 2{,}000$): 전체 11종 (초)}
\scriptsize
\resizebox{\textwidth}{!}{%
\begin{tabular}{@{}lrrrrr@{}}
\toprule
\textbf{기법} & \textbf{BA} ($d\!=\!5$) & \textbf{ER} ($d\!=\!6$) & \textbf{WS} ($d\!=\!8$) & \textbf{Grid} ($d\!=\!88$) & \textbf{PLC} ($d\!=\!5$) \\
\midrule
GPU-Dense & 0.030 & 0.029 & 0.029 & \textbf{0.086} & 0.051 \\
GPU-Sparse & \textbf{0.031} & \textbf{0.025} & \textbf{0.027} & 0.204 & \textbf{0.022} \\
D-STORM-Sparse & 0.395 & 0.379 & 0.380 & 3.371 & 0.326 \\
D-STORM-Dense & 0.329 & 0.325 & 0.417 & 3.906 & 0.282 \\
SciPy (C) & 0.300 & 0.305 & 0.277 & 0.150 & 0.297 \\
I-AORM & 0.459 & 0.479 & 0.584 & 5.630 & 0.430 \\
M-AORM & 0.483 & 0.514 & 0.697 & 6.573 & 0.428 \\
GB-frontier & 0.476 & 0.481 & 0.492 & 0.631 & 0.474 \\
GB-bfs & 1.162 & 1.215 & 1.294 & 3.646 & 1.157 \\
NetworkX & 0.967 & 1.147 & 1.045 & 0.847 & 0.983 \\
\bottomrule
\end{tabular}}
\end{table}

\textbf{발견 9c:}
\begin{itemize}
  \item \textbf{소세계 그래프} (BA, ER, WS, PLC: $d \leq 8$): GPU-Sparse가 최고. CPU 내에서는 D-STORM $>$ SciPy $>$ GB-frontier.
  \item \textbf{Grid} ($d\!=\!88$): \textbf{GPU-Dense가 유일한 0.1초 이하} (0.086s). SciPy(0.150s)가 CPU 최고. D-STORM은 3--4초로 이 위상에서 약함 --- 88회 반복의 Python/Cython 루프 오버헤드.
  \item \textbf{GB-frontier}는 Grid에서 D-STORM보다 5.3배 빠름 (0.631 vs 3.371) --- GraphBLAS C 커널의 반복 효율.
  \item \textbf{I-AORM, M-AORM}은 모든 위상에서 최하위권 --- dense $O(n^3)$ matmul의 한계.
\end{itemize}

\subsection{핵심 발견 종합}

\begin{tcolorbox}[colback=green!5, colframe=green!50!black, title=Exp 9: 전체 통합 벤치마크 핵심 결론]
\begin{enumerate}
  \item \textbf{GPU 가속이 모든 조건에서 최고 성능.} Facebook에서 GPU-Dense 0.126s (CPU 대비 10--100배). BA $n\!=\!5{,}000$에서 GPU-Sparse 0.094s (CPU 대비 21배).
  \item \textbf{CPU 내에서 D-STORM이 대부분의 조건에서 최고.} SciPy(C)와 거의 동등하거나 우위 (소세계). GB-frontier보다 1.2--1.4배 빠름.
  \item \textbf{예외: 고지름 그래프.} Grid($d\!=\!88$)에서 SciPy $>$ GB-frontier $>$ D-STORM. GPU-Dense만이 0.1초 이하 달성.
  \item \textbf{AORM 원본(I/M)은 모든 조건에서 D-STORM에 열위.} I-AORM은 D-STORM-Sparse보다 1.3--5배 느림.
  \item \textbf{GB-bfs(소스별 BFS)는 Python 루프로 인해 비경쟁적.} C/C++ 수준 LAGraph 비교가 추후 필요.
\end{enumerate}
\end{tcolorbox}

\section{종합}

\begin{table}[ht]\centering
\caption{D-STORM 강점과 한계 (GraphBLAS + GPU 비교 포함)}
\begin{tabular}{@{}lp{9.5cm}@{}}
\toprule
\multicolumn{2}{@{}l}{\textbf{강점}} \\
\midrule
S1 & Facebook CPU: I-AORM 대비 3.0배, SciPy(C) 대비 1.2배, NX 대비 8.8배, GB-frontier 대비 1.2배 \\
S2 & Facebook GPU: CPU D-STORM 대비 \textbf{10.1배}, NX 대비 \textbf{85.5배} \\
S3 & 홉 제약 $k=5$: NX 대비 18.3배 (89\% 정확도) \\
S4 & D-STORM 엣지 추가 22--33배 (vs 전체 재계산) \\
S5 & GPU speedup이 그래프 크기에 비례 증가: $n\!=\!5{,}000$에서 \textbf{22.7배} \\
S6 & 소세계 그래프에서 동일 알고리즘의 GraphBLAS 구현보다 1.2--1.4배 빠름 (CPU) \\
\midrule
\multicolumn{2}{@{}l}{\textbf{한계}} \\
\midrule
W1 & Grid($d=88$): SciPy가 CPU 최고, GB-frontier가 CPU D-STORM보다 5.6배 빠름 \\
W2 & 엣지 삭제: 전체 재계산 (rebuild heuristic) \\
W4 & GPU dense 경로는 $n > 10{,}000$에서 VRAM 한계 \\
\bottomrule
\end{tabular}
\end{table}

\end{document}
